No projeto de pesquisa que submeti ao edital de mestrado supracitado, articulei o antropoceno à tese do excepcionalismo humano. Segundo esta, seres humanos são uma exceção no mundo natural, sendo capazes de compreendê-lo e manipulá-lo segundo sua vontade de modos como nenhum outro ser natural poderia. Naquele momento, minha pesquisa deixava de se ocupar com os problemas para se falar do autismo, e começava a se ocupar com os problemas de falar do antropoceno e de viver nele.
Entendo que falar de \enquote{humanismo} é mais assertivo do que falar de \enquote{excepcionalismo humano}. Entendo que o humanismo assume o excepcionalismo humano e nomeia, ao mesmo tempo, um pensamento e um processo de moldura contínua do que é ser humano conforme o conhecimento que se assume acumular por uma humanidade una. Ao nomear o antropoceno, os humanos se auto-nomeiam mais assertivamente do que o fazem à época ou processo que pretendem nomear; a nomeação revela todo conhecimento humanista como um decalque do que é o humano segundo sua cosmologia; ironicamente, a dificuldade em estabelecer um outro referente, no mundo, que sustente a nomeação, revela também os limites desta cosmologia, e demarca sua crise derradeira.
Optei por falar do antropoceno pois estes mesmos limites são aqueles que encontrei ao tentar falar do autismo -- ali também se tratava de uma controvérsia de nomeação para a qual não se conseguia encontrar um referente comum no mundo -- e assim poderia me articular com mais interlocutores. Encontrei no humanismo um problema em comum entre as falas e ações perante o autismo, e as falas e ações perante o antropoceno, e segui pelo caminho mais interessante para falar do problema. Acredito que também estou fazendo ciência orientada pela e para a comunicação quando busco articular o meu problema ao problema do outro e decido tratar e falar deste, e não daquele.